\subsection{Completion sets of pruning supports}
\label{app:pruning-contexts}

Fix one owner and its family $\mathcal F$ of legal supports. For any
family $A$ of supports put
\[
 A^\perp=\{U\subseteq W:\ S\cup U\in\mathcal F\text{ for every }S\in A\}.
\]
The members of $A^\perp$ are the extra labels that make every support
in $A$ legal. Write $A*B=\{S\cup T:S\in A,\ T\in B\}$. At a gate
of the fixed owner the support families multiply by $*$, and at
a foreign gate they take ordinary union.

\begin{theorem}[The local context quotient]
\label{core:thm:pruning-contexts}
Two words are interchangeable in all one-hole contexts testing
legality at the fixed owner exactly when their pruning families
have equal completion sets. It suffices to test the hole alone
and gates at that owner adjoining a fold of extra inputs.
\end{theorem}
\begin{proof}
At a foreign choice the completion sets intersect. At an owned choice,
$U\in(A*B)^\perp$ exactly when $T\cup U\in A^\perp$ for every
$T\in B$. By commutativity, the resulting completion set depends
only on $A^\perp$ and $B^\perp$. Since a family is legal exactly
when its completion set contains $\varnothing$, induction gives
interchangeability.
Conversely, a nonempty distinguishing completion $U$ can be adjoined
by an owned gate whose other child folds all inputs in $U$. This child
has sole support $U$. The hole itself tests the empty completion.
\end{proof}

Equivalently the local states are double polars $A^{\perp\perp}$ of
pruning families. Double polarity gives a closure, since $S\cup U$
is symmetric in $S,U$. Its closed families form a lattice with meet
intersection and join closure of union. The owned gate is multiplication
followed by closure. The product law in the proof makes multiplication
associative on closed families and distributive over joins.
If $\mathcal F$ is union closed, each polar is union closed, and
multiplication is idempotent. A state is legal exactly when it is
contained in the closed family $\{\varnothing\}^\perp=\mathcal F$.
Thus one can compute joint feasibility in the product of these
intrinsically specified local algebras. At each gate exactly the
owner's coordinate uses multiplication and all other coordinates
use join.

\paragraph{When the local quotients combine independently.}
Call a word universally extendible if all its completion sets
contain every subset of $W$. For constraints given by implications
$j\Rightarrow\operatorname{hit}(T)$ with $j\notin T$, this means
every pruning at an owner meets every tail $T$ at that owner.
Indeed a support missing $T$ can be completed by adjoining $j$,
which violates the implication. Such a word therefore exists
exactly when the singleton profile demands for all the tails meet
the common-intersection condition of \cref{core:thm:profiles}.
Use $W$ as a redundant demand at each owner.

\begin{theorem}[Independence of the local quotients]
\label{core:thm:pruning-context-product}
Suppose $K\ge3$ and a universally extendible word exists. Every
tuple of separately attainable local completion sets is jointly
attainable. Equality in this product is exactly equivalence in
all one-hole contexts testing legality at every owner together.
If every tail is nonempty, full product attainability conversely
implies the common-intersection condition on the tails.
\end{theorem}
\begin{proof}
Attach a universally extendible word by a gate at owner $s$.
This replaces the completion set at $s$ by the entire powerset,
and leaves the other completion sets unchanged. For an individual
local witness, do this at all other owners. To combine the resulting
witnesses, the derived union of \cref{game:thm:union} takes intersection
of completion sets at every coordinate. In its two inner gates
and one outer gate, an owner outside the three sees only unions.
An inner owner sees $A\cup B$ together with $A*B$, and an outer
owner sees $(A\cup B)*(A\cup B)$. Since legality is closed under
unions, the completion sets in all these cases are precisely
$(A\cup B)^\perp$. This gives a common word for the tuple.

Local equivalence implies joint equivalence. To distinguish a
local discrepancy, first apply its distinguishing context, then
attach the universally extendible word at every other owner.
The joint acceptance test now reads that one local coordinate.
Finally if all tails are nonempty, each owner's local top is
attainable by its fold on all inputs. Full product attainability
must produce the tuple of local tops, giving the universal word.
\end{proof}

Joint equivalence can be strictly coarser even if one feasible word
uses every input. On $\{a,b,c\}$ impose $a\Rightarrow b$ at owner
$r$ and $b\Rightarrow c$ at owner $s$. The word
$(a\star_r b)\star_s c$ is feasible. The projection $a$ and the word
$a\star_s c$ have the same violation at $r$ and differing completion
sets at $s$. In any surrounding context their $r$-violation persists
until the first $r$-gate whose sibling has all $r$-supports meeting
$\{b\}$. Such a sibling has the $s$-pruning $\{b\}$ by the foreign
singleton argument of \cref{core:prop:compulsory-prunings}.
At that gate it erases the $r$-violation and introduces the same
$s$-violation in both words. Their two states then agree for the
rest of the context. If there is no such gate, both words are rejected.

\paragraph{Size of a local quotient.}
An implication $j\Rightarrow\operatorname{hit}(T)$ contributes
one of three conditions on an extra set $U$. If every support of the
word hits $T$, the row imposes no condition. If some support misses
$T$ but none of those contains $j$, the extra set must obey the
original implication. If some support containing $j$ misses $T$,
the extra set must meet $T$. Thus $m$ implications at an owner give
at most $3^m$ context classes.

\begin{proposition}[Exact context count at one head]
\label{core:prop:pruning-context-index}
Suppose an owner's only constraints are at one head $j$, with an
antichain of $m$ nonempty proper tails in the passive input set.
With a second gate owner available, there are exactly $3^m$
context classes. Exactly $2^m$ occur among words that themselves
pass legality.
\end{proposition}
\begin{proof}
For a tail $T$, its complement in the passive inputs meets every
other tail and misses $T$. Use that complement and the same set
with $j$ adjoined as test completions. They distinguish the three
conditions at $T$ and impose no restriction from the other rows.
To realize any assignment of conditions, include the fold on all
passive inputs. For each tail requiring the original implication,
include a fold on its passive complement. For each tail requiring
a hit, include that complement with $j$ adjoined instead. Take
the union of these support families by gates at a foreign owner.
Each added support changes exactly the specified condition. The
empty completion is legal precisely when none of the conditions
requires a hit.
\end{proof}

For example, let the passive inputs be $d\ge2$ disjoint pairs.
Take as tails all choices of one input from each pair. At the
head the positive prime guards exclude both postings in any
single pair, so there are just $d$ primes on each relevant facet.
The local context quotient nevertheless has $3^{2^d}$ classes,
including $2^{2^d}$ classes of already legal words. Thus even a
single owner's omissions can require doubly exponentially many
states in a deterministic compositional summary on the primes.

\begin{proposition}[Short distinguishing contexts]
\label{core:prop:pruning-context-length}
For $m$ listed pure rows, every one-hole context has a subcontext
with at most $3m$ gates on the path to the hole and the same effect
on all row statuses for every filling of the hole. Deciding joint
context equivalence on listed pure rows is PSPACE-complete at
every fixed $K\ge3$, already for two projections.
\end{proposition}
\begin{proof}
Number the three row statuses $0,1,2$ in the displayed order.
On each row a foreign gate on the path from the hole takes maximum
with the corresponding sibling status. An owned gate with sibling
$0$ resets to $0$, and an owned gate with nonzero sibling takes
maximum unless the incoming status is $0$. Thus the unary action
of a gate is either constant or pointwise extensive. For one row
keep the last constant action, if any, and the later actions that
change the accumulated map. If no constant occurs, start at the
identity $(0,1,2)$, whose coordinate sum can increase at most three
times. After a constant action, every subsequent map is itself
constant, so at most two increases are possible. Keep the union of
these selected gates over all rows. For any row the retained actions
before its last reset are erased by that reset. Subsequently each
discarded action was inert, and each retained action has the same
input map as originally, by induction. This gives the bound.

To find a distinguishing context, guess its path and the exact row
statuses of each sibling. A sibling status vector is realizable
exactly when its three possible values on each pure probe are jointly
realizable, a subpower membership test in polynomial space by
\cref{core:thm:subpower-complexity}. These tests and the path
calculation use polynomial space. Equivalence has the same upper
bound by closure under complement.

For hardness use the PSPACE-complete problem of omitting the listed
monochromatic rows and supplying one specified source $j$ at the
empty coloring (\cref{core:thm:two-color-pspace}). This asks for an
accepted word containing $j$. Add two fresh inputs and require that
they occur together in every pruning at each of two distinct owners.
The pairing exclusion argument of \cref{core:thm:prime-pruning-hardness}
shows that neither can occur in an accepted word. Compare $j$ with
one fresh projection. A context accepting the latter is impossible.
A context accepting $j$ is exactly a witness to the original
problem, with a particular occurrence of $j$ distinguished as the
hole. Thus distinguishability, and hence also its complement, is
PSPACE-hard.
\end{proof}

\begin{corollary}[Exact interpolation with jointly attainable hits]
\label{core:cor:pruning-context-interpolation}
Consider a list of pure rows with a prescribed status at each row:
all prunings hit its tail, some pruning misses it but none containing
the head does, or some pruning containing the head misses it. Suppose
the tails of rows with one of the first two statuses satisfy the
common-intersection condition. Then simultaneous realization is
equivalent to separate realization at each owner. Given the rows
explicitly, both checks take polynomial time at fixed $K\ge3$.
\end{corollary}
\begin{proof}
A foreign choice takes the maximum in the displayed order. An owned
choice makes the all-hit status absorbing and otherwise takes the
maximum. The assumed intersections provide a word hitting all tails
with one of the first two statuses. Adjoin it at the foreign owners
of each local witness. This forces all-hit at their constrained
rows. At the selected owner each upper-bounded status stays the
same. A row with the third status stays there since foreign choice
takes maximum. The derived union takes maximum at every owner, so
it combines the padded witnesses.

For a local test put the upper bounds on each support. Require it
to meet every tail labeled all-hit, and to satisfy the displayed
head implication for every row with the second status. Each row
without the all-hit status additionally needs a witness support
inside $W\setminus T$. For the third status this support must contain
the head. To test a witness, start with $W\setminus T$ and repeatedly
delete a head whose required tail misses the current set. The
remaining set is greatest among those obeying the conditional
requirements in that domain. It succeeds if nonempty, if it meets
all unconditional requirements, and if it contains the head when
needed. With no witness rows, test on $W$ instead. Take the fold
of each resulting support and combine them at a foreign gate owner.
This realizes the local statuses. The loop takes at most $|W|$
deletions per row and the common-intersection test examines at most
one tail per owner.
\end{proof}
